\chapter{The Alarm Signal}
% full version: chapters/14_pain.tex

\pencilfig{14}{\pencilart{14}}{Pain is a fire alarm, not a thermometer.}

All the other senses report what the world is like. Pain does not. It reports that something threatens the body, and its loudness is set largely by the machine itself, not by the stimulus. Pain is a fire alarm, not a thermometer.

Pain sensors fire only at something strong: at heat above forty-odd degrees Celsius, at heavy pressure, at caustic substances. They adapt much less than touch sensors do: an alarm must not fall silent while the threat is still there.

\section*{Two Wires}

Remember hitting your finger. First comes a sharp, precise pain: where and when. A moment later, a dull, spread-out, lasting one: how bad. These are two different wires. One is fast; its signal arrives in tens of milliseconds. The other is slow; its signal takes on the order of a second.

\section*{The Gate}

Why do we rub a bruise after bumping it? In the spinal cord the pain signal passes through a ``gate.'' Touch wakes a braking neuron that half-closes it. So rubbing a bruise really does help, but only against moderate pain; strong pain gets through anyway. That strong pain itself switches off the braking neuron is part of the original gate theory, and there is no direct confirmation of it yet.

\pencilfig{126}{%
% gate: the pain wire drives the relay neuron; touch wakes an inhibitory neuron that closes the gate
\draw[pen=pcAmber] (0,1.35) -- (3.05,1.0); \draw[arr=pcAmber] (2.8,1.03) -- (3.08,1.0);
\node[lbl, text=pcAmber!70!black, anchor=east] at (-0.05,1.35) {pain};
\draw[pen=pcBlue] (0,-0.15) -- (1.75,-0.15); \draw[arr=pcBlue] (1.5,-0.15) -- (1.78,-0.15);
\node[lbl, text=pcBlue, anchor=east] at (-0.05,-0.15) {touch};
\draw[dotpen=pcRed, wash=pcRed] (2.05,-0.15) circle (0.26);
\draw[pcRed, line width=0.7pt, -{Bar[width=6pt]}] (2.25,0.05) -- (3.2,0.66);
\node[lbl, text=pcRed, anchor=north] at (2.05,-0.45) {brake};
\draw[dotpen=pcGray, wash=pcGray] (3.4,0.9) circle (0.34);
\node[lbl, anchor=south] at (3.4,1.3) {gate};
\draw[pen] (3.75,0.9) -- (5.6,0.9); \draw[arr] (5.35,0.9) -- (5.65,0.9);
\node[lbl, anchor=west] at (5.7,0.9) {to the brain};
}{The gate in the spinal cord. The pain signal goes to the brain through a gate neuron, and touch wakes a braking neuron that half-closes the gate. So rubbing a bruise helps, but only against moderate pain.}

\section*{A Volume Knob from Above}

The brain can make pain louder or quieter by itself. Radio channels run down to the gate from the brain above: \nt{SERO}, \nt{NORA}, and special substances resembling painkillers that the brain makes itself. In a dangerous situation pain can almost vanish: someone running from danger has no time to limp. Expecting relief also dulls pain, partly through the same substances.

\section*{When the Alarm Gets Stuck}

Repeated pain raises sensitivity: the same blow starts to hurt more. This is learning too: a loop that turns itself up while the wound heals. But a strong loop has a danger: it can ``stick'' and keep ringing after the wound has healed, like the group of neurons from the second chapter that keeps itself going.

\section*{Teacher and Interrupt}

Pain is also a teacher: for dopamine's ``better or worse than expected,'' it works as a negative reward. And it interrupts what you are doing, like a burst of noradrenaline: drop everything, deal with the danger. And the fastest reaction, pulling your hand away from something hot, closes right in the spinal cord through that same pair of opposing muscles from the last chapter, before you have even had time to feel the pain.

\fullversion{14}
